\documentclass[a4paper]{memoir}
\usepackage{graphicx}
%% Language and font encodings
\usepackage{microtype}
\usepackage[paperheight=88.9mm,paperwidth=63.5mm,margin=5mm]{geometry}
\usepackage{ragged2e}
\usepackage{adjustbox}
\usepackage{pgffor}
\usepackage{emoji}
\usepackage{tikz}
\newcommand{\cardtext}{0-$\infty$}
\newcommand{\grouptext}{2-10}
\newcommand{\agetext}{12+}
\newcommand{\clocktext}{5$''$}
\newcommand{\languagecrop}{french_crop}
\newcommand{\colorcount}{6}
\newcommand{\cardtitle}{PR\'ESIDENT}
\newcommand{\cardsubtitle}{FANT\^OME}
\def\colors{0/myyellow, 1/myorange, 2/myred, 3/mypurple, 4/myblue, 5/mygreen}
\begin{document}
{\footnotesize
\input{../../../rule_lib/rule_lib.tex}

\noindent
\textbf{\emoji{arrows-counterclockwise} Mise en place.}  
Distribue les cartes en incluant un joueur fictif. Ses cartes sont écartées. Le joueur avec le \textbf{10 jaune} commence.
\\
\\
\noindent
\textbf{\emoji{dart} But du jeu.}  
Être le premier à poser toutes ses cartes.
\\
\\
\noindent
\textbf{\emoji{video-game} Déroulement.}  
À ton tour, tu peux :
\begin{itemize}
    \item \textbf{Poser une combinaison} de cartes identiques (singleton, paire, triple...).  
    Tu dois jouer le même nombre de cartes et une valeur \textbf{supérieure ou égale} à la combinaison précédente.
    
    \item \textbf{Passer ton tour.}  
    Dans ce cas, tu ne pourras plus rejouer tant que le cycle en cours n’est pas terminé.
\end{itemize}

\noindent
\emoji{pushpin} Les \textbf{jokers} peuvent remplacer n’importe quelle carte dans une combinaison.
\\
\\
\noindent
\emoji{warning} Le cycle s’arrête si :
\begin{itemize}
    \item Tous les autres joueurs passent.
    \item Une combinaison de 12 est jouée.
    \item Une combinaison de même valeur que la précédente, d’au moins 2 cartes, sans joker, est jouée.
\end{itemize}
Le dernier joueur à avoir joué  commence le cycle suivant.
\\
\\
\emoji{pushpin} \textbf{Interdit de finir sur une combinaison de 12} — sinon, tu perds.

}
\end{document}



